\documentclass[../book.tex]{subfiles}
\begin{document}

\section{重要公式}

\subsection{三角函数}

\subsubsection{诱导公式}

奇变偶不变，符号看象限（\textbf{\underline{等式左边的角所在象限的正负}}）

\begin{enumerate}
    \item $\sin (\displaystyle\frac{\pi}{2} \pm \alpha) = \cos \alpha$
    \item $\sin (\pi \pm \alpha) = \mp \sin \alpha$
    \item $\cos (\displaystyle\frac{\pi}{2} \pm \alpha) = \mp \sin \alpha$
    \item $\cos (\pi \pm \alpha) = \mp \cos \alpha$
\end{enumerate}

\subsubsection{和差公式}

\begin{enumerate}
    \item $\sin (\alpha \pm \beta) = \sin \alpha \cos \beta \pm \cos \alpha \sin \beta$
    \item $\cos (\alpha \pm \beta) = \cos \alpha \cos \beta \mp \sin \alpha \sin \beta$
    \item $\tan (\alpha \pm \beta) = \displaystyle\frac{\tan \alpha \pm \tan \beta}{1 \mp \tan \alpha \tan \beta}$
    \begin{itemize}
        \item $\tan (\displaystyle\frac{\pi}{4} - \alpha) = \displaystyle\frac{1 - \tan \alpha}{1 + \tan \alpha}$
    \end{itemize}
    \item $\cot (\alpha \pm \beta) = \displaystyle\frac{\cot \alpha \cot \beta \mp 1}{\cot \alpha \pm \cot \beta}$
\end{enumerate}

\subsubsection{倍角公式}

\begin{enumerate}
    \item $\sin 2\alpha = 2\sin \alpha \cos \alpha$
    \item $\cos 2\alpha = \cos^2 \alpha - \sin^2 \alpha = 1 - 2\sin^2 \alpha = 2\cos^2 \alpha - 1$
    \item $\sin 3\alpha = -4\sin^3 \alpha + 3\sin \alpha$
    \item $\cos 3\alpha = 4\cos^3 \alpha - 3\cos \alpha$
\end{enumerate}

\subsubsection{半角公式}

\begin{enumerate}
    \item $\sin^2 \displaystyle\frac{\alpha}{2} = \displaystyle\frac{1}{2}(1 - \cos \alpha)$
    \item $\cos^2 \displaystyle\frac{\alpha}{2} = \displaystyle\frac{1}{2}(1 + \cos \alpha)$
\end{enumerate}

\subsection{一元二次方程}

一元二次方程 $ax^2 + bx + c = 0(a \neq 0)$

\begin{enumerate}
    \item 判别式$\Delta = b^2 - 4ac$
    \begin{itemize}
        \item $\Delta \geq 0$，方程有两个实根$x_{1, 2} = \displaystyle\frac{-b \pm \sqrt {b^2 - 4ac}}{2a}$
        \item $\Delta < 0$，方程有两个共轭的复根$x_{1, 2} = \displaystyle\frac{-b \pm \sqrt {4ac - b^2 i}}{2a}$
    \end{itemize}
    \item 根与系数的关系（韦达定理）
    \begin{itemize}
        \item $x_1 + x_2 = -\displaystyle\frac{b}{a}$
        \item $x_{1}x_2 = \displaystyle\frac{c}{a}$
    \end{itemize}
\end{enumerate}

\subsection{因式分解}

\begin{enumerate}
    \item $(a + b)^2 = a^2 + 2ab + b^2$
    \item $(a - b)^2 = a^2 - 2ab + b^2$
    \item $(a + b)^3 = a^3 + 3a^{2}b + 3ab^2 + b^3$
    \item $(a - b)^3 = a^3 - 3a^{2}b + 3ab^2 - b^3$
    \item $a^2 - b^2 = (a + b)(a - b)$
    \item $a^3 - b^3 = (a - b)(a^2 + ab + b^2)$
    \item $a^3 + b^3 = (a + b)(a^2 - ab + b^2)$
    \item 二项式定理
    \[
        \begin{aligned}
        (a+b)
            ^n
            &= \sum_{k=0}^{n} C_n^k a^{\,n-k} b^{\,k} \\
            &= a^n
            + n a^{n-1} b
            + \frac{n(n-1)}{2!} a^{n-2} b^2 \\
            &\quad + \cdots
            + \frac{n(n-1)\cdots(n-k+1)}{k!} a^{n-k} b^k \\
            &\quad + \cdots
            + n a b^{n-1}
            + b^n
        \end{aligned}
    \]
\end{enumerate}

\subsection{阶乘与双阶乘}

\begin{enumerate}
    \item $n! = 1\cdot 2\cdot 3\cdots \cdot n$，规定$0! = 1$
    \item $(2n)!! = 2\cdot 4\cdot 6\cdots \cdot 2n = 2^n\cdot n! $
    \item $(2n - 1)!! = 1\cdot 3\cdot 5\cdots \cdot (2n - 1) $
\end{enumerate}

\end{document}
